\documentclass[11pt,letterpaper]{article}
\usepackage{cornell-notes}
\usepackage{mathtools}

\title{Chapter 13: Extending the Core Problems}
\author{}
\date{}

\setDocTitle{Chapter 13: Extending the Core Problems}
\setDocSubtitle{String Algorithms}
\setDocVersion{v1.0}
\setDocStatus{Study Companion}
\setDocOwner{String Algorithms Cornell Notes}
\setDocAuthor{}
\setDocDate{\today}
\setDocSummary{Cornell-format study and review notes for Chapter 13: Extending the Core Problems.}

\setCornellCollection{String Algorithms}
\setCornellUnitType{Chapter}
\setCornellUnitNumber{13}
\setCornellUnitTitle{Extending the Core Problems}
\setCornellCompanionLabel{Study and review companion}

\hypersetup{pdftitle={Chapter 13: Extending the Core Problems},pdfsubject={Cornell notes},pdfauthor={}}

\begin{document}
\maketitle

\begin{CornellOverviewBox}{Chapter overview}
\textbf{Purpose.} Move beyond one optimum under fixed weights by analyzing parameter regions, near-optimal solutions, and chains of local alignments.

\textbf{Learning objectives}
\begin{itemize}
  \item Interpret an alignment score as a linear function of parameters.
  \item Describe the polygonal decomposition of parameter space.
  \item Count or enumerate near-optimal alignments without confusing paths and scores.
  \item Optimize a compatible chain of local alignments.
\end{itemize}

\textbf{Prerequisites.} Alignment DP, edit graphs, affine functions, convex geometry, DAG paths, and interval scheduling.
\end{CornellOverviewBox}

\begin{CornellNotesTable}
\CornellNoteRow{\S13.1: Why make alignment parametric?}{A fixed mismatch or gap penalty can select one alignment while a nearby parameter choice selects another. Parametric alignment asks for all regions of parameter space in which the same alignment summary is optimal, exposing sensitivity rather than hiding it behind one arbitrary setting.}
\CornellNoteRow{How is a score represented?}{For an alignment $A$, collect counts such as matches, substitutions, and gap events into a feature vector $f(A)$. With parameter vector $\theta$, the score is $\theta\cdot f(A)$, an affine function. The optimum over alignments is the upper or lower envelope of these functions.}
\CornellNoteRow{What is the polygonal decomposition?}{Parameter values yielding the same optimal feature vector form a convex region. Boundaries occur where two alignment-score functions tie. In two dimensions these regions form a polygonal subdivision; in higher dimensions they form a polyhedral complex.}
\CornellNoteRow{How are regions computed?}{Use an alignment oracle for selected parameter values, add the resulting feature vectors to a convex hull or envelope, and recurse where current supporting solutions do not certify a region. Efficiency depends on the number of distinct optimal summaries, not the number of alignment paths.}
\CornellNoteRow{What can parametric output reveal?}{Stable regions indicate robust conclusions; narrow regions indicate sensitivity. The decomposition supports parameter selection, comparison of scoring schemes, and identification of breakpoints where an optimal alignment changes.}
\CornellNoteRow{\S13.2: What is a suboptimal alignment?}{Define a tolerance from the optimum, such as additive score loss or a bounded rank. The task may be to count, enumerate, sample, or compactly represent all qualifying paths. These are different outputs with different complexity and duplication concerns.}
\CornellNoteRow{Why reweight the alignment graph?}{Use optimum prefix or suffix potentials to transform edge weights so the cost above optimum is nonnegative and accumulates along deviations. Near-optimal paths then correspond to paths with bounded reweighted cost, making counting or ordered enumeration easier.}
\CornellNoteRow{What causes combinatorial explosion?}{Many distinct traceback paths can have the same score or even the same displayed alignment under zero-cost choices. Enumeration time must include output size, and a practical method needs a policy for ties, duplicates, and a maximum number of returned alternatives.}
\CornellNoteRow{\S13.3: What is chaining?}{Local alignments provide candidate anchors with coordinates and scores. A chain selects anchors that are mutually compatible in both sequence orders, possibly charging gap costs between them, to produce a coherent larger-scale comparison.}
\CornellNoteRow{How is the best chain computed?}{Sort anchors by endpoints and use dynamic programming: the best chain ending at anchor $j$ is its score plus the best compatible predecessor contribution. Range-search or geometric data structures accelerate predecessor maxima when the naive all-pairs transition is too costly.}
\CornellNoteRow{What is the chaining invariant?}{Every stored predecessor ends before the current anchor in each required coordinate and therefore cannot cross it. The DP value is optimal among chains ending at that anchor. Coordinate ties and overlapping anchors require an explicit compatibility rule.}
\CornellNoteRow{\S13.4: What should practice emphasize?}{Draw a small parameter-space envelope, distinguish an alignment from its feature vector, reweight a tiny DAG, and trace a chain DP with crossing and noncrossing anchors.}
\end{CornellNotesTable}

\begin{CornellSummaryBox}{Chapter synthesis}
The alignment graph supports more than a single fixed-weight optimum. Parameterization studies which feature summaries win across scoring choices; reweighting exposes the cost of deviations from optimum; chaining assembles local evidence while enforcing coordinate consistency. Output size and equivalence definitions become central.
\end{CornellSummaryBox}

\begin{CornellSummaryBox}{Key takeaways}
\begin{CornellRecallList}
  \item Each alignment defines an affine score function in parameter space.
  \item Optimal feature vectors induce convex parameter regions.
  \item Sensitivity is visible at region boundaries.
  \item Near-optimal path enumeration is inherently output sensitive.
  \item Reweighting turns excess cost into nonnegative deviations.
  \item Chaining requires monotone compatibility in both sequences.
\end{CornellRecallList}
\end{CornellSummaryBox}

\begin{CornellOverviewBox}{Algorithm selection: when to use this}
\begin{CornellChecklist}
  \item Use parametric alignment when the scoring parameters are uncertain or contested.
  \item Use near-optimal enumeration when downstream analysis needs alternatives rather than one traceback.
  \item Use chaining when reliable local anchors should form a larger collinear explanation.
\end{CornellChecklist}
\end{CornellOverviewBox}

\begin{CornellExamBox}{Self-test questions}
\begin{CornellRecallList}
  \item Why is an alignment score affine in its feature counts?
  \item What does a boundary between two parameter regions mean?
  \item Why can many paths map to one feature vector?
  \item How do potentials help enumerate near-optimal paths?
  \item What coordinate conditions make two local alignments chain-compatible?
\end{CornellRecallList}
\end{CornellExamBox}

{\footnotesize
\textbf{Terms and notation to remember.} parameter vector, feature vector, score envelope, polygonal decomposition, optimality region, potential, excess cost, anchor, compatible chain.

\textbf{Connections.} Multiple alignment expands the state dimension; database search and assembly use chaining to connect local or overlap evidence into global structure.
}

\end{document}
